

\lussec 3. Field theory in $\mib D$+1 dimensions

The discussion in the following sections heavily builds 
on the fact that the correlation functions 
of the time-dependent fields coincide with 
the correlation functions in
a local field theory in $D$+1 dimensions,
the extra dimension being the flow time.
The observation dates back to the seminal work of
Zinn--Justin and Zwanziger [\cite{Zinn},\cite{ZinnZwanziger}]
on the Langevin equation and recently
allowed the renormalizability 
of the (pure-gauge) gradient flow to be established 
to all orders of perturbation theory
[\cite{RenFlow}]. 

As will become clear below, the inclusion of the quark fields 
in the theory in $D$+1 dimensions is straightforward.
The general setup is exactly the same as in ref.~[\cite{RenFlow}],
where an introduction to the subject and further details 
can be found. 


\lussubsec 3.1 Fields and action

In addition to the fundamental and the time-dependent fields
already encountered, the theory in $D$+1 dimensions
involves the Lagrange-multiplier fields $L_{\mu}(t,x)$, $\lambda(t,x)$
and $\lambdabar(t,x)$. The latter are fermion fields with
the same indices as the quark fields, while $iL_{\mu}(t,x)$ is a
vector field that takes values in the Lie algebra of $\SUn$. 

As explained in ref.~[\cite{RenFlow}], 
gauge fixing in perturbation theory requires the introduction of 
the Faddeev--Popov ghost
fields $c(x),\cbar(x)$ and the associated 
time-dependent fields $d(t,x),\dbar(t,x)$,
where
\equation{
  \left.d\right|_{t=0}=c.
  \enum
}
No boundary condition is imposed on the field $\dbar(t,x)$, 
which, in many ways, plays a 
r\^ole similar to the Lagrange-multiplier fields.
Apart from the boundary conditions (2.1),(2.6) and (3.1),
all fields are
assumed to be unconstrained at this point.

The total action of the theory,
\equation{
  S_{\rm tot}=S+\SGfl+\SFPfl+\SFfl,
  \enum
}
includes the gauge-fixed QCD action $S$ of the fields at 
flow time zero and the bulk actions
\equation{
  \SGfl=-2\int_0^{\infty}\rmd t\int\rmd^D\kern-1pt x\, 
  \tr\bigl\{L_{\mu}(t,x)
  \bigl(\partial_tB_{\mu}-
        D_{\nu}G_{\nu\mu}-\alpha_0D_{\mu}\partial_{\nu}B_{\nu}
  \bigr)(t,x)\bigr\},
  \enum
  \nexteq{3.0ex}
  \SFPfl=-2\int_{0}^{\infty}\rmd t\int\rmd^D\kern-1pt x\,
  \tr\bigl\{
  \dbar(t,x)\bigl(\partial_td-\alpha_0 D_{\mu}\partial_{\mu}d\bigr)(t,x)
  \bigr\},
  \enum
  \nexteq{3.0ex}    
  \SFfl=\int_0^{\infty}\rmd t\int\rmd^D\kern-1pt x\,
  \bigl\{\lambdabar(t,x)
         (\partial_t-\Delta+\alpha_0\partial_{\nu}B_{\nu})\chi(t,x)
  \noenum
  \nexteq{1.5ex}
  {\phantom{\SFfl=\int_0^{\infty}\rmd t\int\rmd^D\kern-1pt x\,\bigl\{}}
  +\chibar(t,x)
  (\lvec{\partial}_t-\lvec{\Delta}-\alpha_0\partial_{\nu}B_{\nu})
  \lambda(t,x)\bigr\}.
  \enum
}
With respect to the pure-gauge flow studied
in ref.~[\cite{RenFlow}], 
the bulk quark action $\SFfl$ is the only new term in the 
total action.


\lussubsec 3.2 Correlation functions

The $n$-point correlation functions of the 
gauge, quark, ghost and Lagrange-multiplier fields 
are formally defined through the functional integral
over all fields with ``Boltzmann factor'' 
$\exp\{-S_{\rm tot}\}$. 

In the case of the correlation functions of the gauge and quark fields,
the functional integral can be worked out, to some
extent, by performing the integral over the Lagrange-multiplier
fields. Since the action is linear in these fields,
the integration yields 
a product of functional Dirac $\delta$-functions.
The bulk actions were chosen such that
the $\delta$-functions amount to imposing the flow equations on 
the time-dependent gauge and quark fields. 
Up to mathematical subtleties and possibly non-trivial Jacobian factors,
the correlation functions thus coincide with the ones
calculated directly, as in sect.~2, by
solving the flow equations with the initial
conditions (2.1),(2.6) and treating the calculated
time-dependent fields as QCD observables.

The equivalence of the theory in $D$+1 dimensions and the 
direct computation of the correlation functions
can be rigorously shown in perturbation theory
[\cite{RenFlow}] and in lattice QCD (sect.~5).
In particular, all correlation functions 
that do not involve the time-dependent quark fields or
the associated Lagrange-multiplier fields
coincide with the ones studied in ref.~[\cite{RenFlow}],
where the quark fields were not evolved in flow time.
This ``reduction property'' 
of the theory in $D$+1 dimensions strongly constrains
the form of the possible counterterms
required for renormalization (see subsect.~3.4).


\lussubsec 3.3 Fermion integral

Since the action in $D$+1 dimensions 
is quadratic in the fermion fields,
the functional integral over these fields can be
evaluated (at fixed bosonic fields) by applying Wick's theorem.
The basic Wick contractions are worth being given explicitly, 
because the expressions provide some insight
into the structure of the theory
in $D$+1 dimensions and will, in any case, be
frequently needed later.

The Wick contractions of the basic fields
are determined by the action and the boundary conditions (2.6).
In particular, the propagator of
the quark fields at flow time zero,
\equation{
  \wick{\psi(x)\psibar(y)}=S(x,y),
  \enum
  \nexteq{2.5ex}
  (\slash{D}+M_0)S(x,y)=\delta(x-y),
  \enum
}
coincides with the one in $D$ dimensions.
Apart from the quark propagator $S(x,y)$, 
the contractions of the time-dependent fermion fields,
\equation{
  \longwick{\lambda(t,x)\lambdabar(s,y)}=0,
  \enum
  \nexteq{2.5ex}
  \longwick{\chi(t,x)\lambdabar(s,y)}=\theta(t-s)K(t,x;s,y),
  \enum
  \nexteq{2.5ex}
  \longwick{\lambda(t,x)\chibar(s,y)}=\theta(s-t)K(s,y;t,x)^{\dagger},
  \enum
  \nexteq{2.5ex}
  \longwick{\chi(t,x)\chibar(s,y)}=\int\rmd^Dv\,
  \rmd^Dw\,
  K(t,x;0,v)S(v,w)K(s,y;0,w)^{\dagger},
  \enum
}
involve the fundamental solution
\equation{
  \{\partial_t-\Delta+
  \alpha_0\partial_{\nu}B_{\nu}\}K(t,x;s,y)=0\quad\hbox{if}\quad t\geq s,
  \enum
  \nexteq{2.0ex}
  \lim_{t\to s}K(t,x;s,y)=\delta(x-y),
  \enum
}
of the quark flow equation. Note that the
kernel $K(t,x;s,y)$ is flavour-independent, proportional
to the unit Dirac matrix and 
an $N\times N$ complex matrix in colour space (the adjoint of
the kernel in eqs.~(3.10) and (3.11) is 
taken in index space only). 

Consistently with the boundary conditions at vanishing flow time,
the contractions 
involving the fundamental quark and antiquark fields,
\equation{
  \longwick{\chi(t,x)\psibar(y)}=\int\rmd^Dv\,K(t,x;0,v)
  S(v,y),
  \enum
  \nexteq{2.5ex}
  \wick{\psi(x)\chibar(s,y)}=\int\rmd^Dw\,S(x,w)
  K(s,y;0,w)^{\dagger},
  \enum
  \nexteq{2.5ex}
  \longwick{\lambda(t,x)\psibar(y)}=
  \wick{\psi(x)\lambdabar(s,y)}=0,
  \enum
}
are included in eqs.~(3.9)--(3.11) as special cases.
Note, however, that the behaviour of 
correlation functions 
near space-time boundaries tends to 
be complicated by short-distance singularities
[\cite{Symanzik}]. 
The limits
\equation{
  \lim_{s\to0}\longwick{\lambda(0,x)\chibar(s,y)}=
  \lim_{t\to0}\longwick{\chi(t,x)\lambdabar(0,y)}=
  \delta(x-y),
  \enum
}
for example,
differ from the contractions (3.16) at $t=s=0$ by a contact term.


\lussubsec 3.4 Renormalization

In perturbation theory, the renormalization properties of
the theory in $D$+1 dimensions can
be determined using power counting and 
symmetry arguments.
The task of figuring out all possible local counterterms
to the action is simplified by the fact that the renormalization of
the theory where only the gauge field is evolved in flow time
is already known to require no more than the usual
QCD parameter and field renormalizations [\cite{RenFlow}]. In particular,
at flow time $t>0$,
the fields $B_{\mu}(t,x)$ and $L_{\mu}(t,x)$ need not
be renormalized. 

When the time-dependent quark fields are included in the theory,
further counter\-terms may need to be added to the total action.
The reduction property mentioned at the end of subsect.~3.2
however implies that any such term
must be proportional to the quark fields or
the associated Lagrange-multiplier fields.
Moreover,
as explained in subsect.~7.2 of ref.~[\cite{RenFlow}], 
bulk counterterms (i.e.~local terms integrated over all
$D$+1 coordinates)
are excluded, because
the correlation functions at large flow times
are given by tree diagrams. 

Taking the symmetries of the unrenormalized theory 
into account, and the fact
that the quark and Lagrange-multiplier fields have
power-counting dimension $3/2$ and $5/2$, respectively,
a moment of thought reveals that the only possible additional
counterterm is proportional to
\equation{
  \int\rmd^D\kern-1pt x\,
  \left\{\lambdabar(0,x)\psi(x)+\psibar(x)\lambda(0,x)\right\}.
  \enum
}  
The insertion of this term in the Feynman diagrams
is equivalent to a renormalization 
\equation{
  \chi=Z_{\chi}^{-1/2}\ren{\chi},
  \qquad
  \lambda=Z_{\chi}^{1/2}\ren{\lambda},
  \enum
}
of the fermion and, similarly, the anti-fermion fields
at positive flow time.
Apart from the usual renormalization of the 
QCD parameters and fields,
the renormalization of the theory in $D$+1
dimensions thus requires the time-dependent quark
and associated Lagrange multiplier fields 
to be renormalized as specified above.
